\begin{ejercicio}
    % ID: 1
    Mostrar que $f_*$ está bien definida y satisface que
    \begin{enumerate}[label=(\arabic*)]
        \item  ${\id_X}_*=\id_{\pi_0(X)}.$
        \item  Para cualquier $f:X\to Y$ y $g:Y\to Z$ funciones continuas, entonces $(g\circ f)_*=g_*\circ f_*$.
    \end{enumerate}
%\begin{dem}
%    Supongamos $x\simeq y$, entonces existe una trayectoría continua $\sigma:I\to X$ tal que $\sigma(0)=x$ y $\sigma(1)=y$. Entonces $f\circ \sigma(0)=f(x)$ y $f\circ \sigma(1)=f(y)$
%\end{dem}
\end{ejercicio}

\begin{ejercicio}
    % ID: 2
    Mostrar que la relación de homotopía en el conjunto $\Omega(X,x_0)$ es una relación de equivalencia.
    %\begin{dem}
    %   Supongamos
    %\end{dem}
\end{ejercicio}

\begin{ejercicio}
    % ID: 3
        Mostrar que la homotopía es una relación de equivalencia en el conjunto de funciones continuas de $X$ en $Y$, denotado por $\text{Map}(X,Y)$.
    \end{ejercicio}

\begin{ejercicio}
    % ID: 4
        Muestra que si $X\cong Y$ (homeomorfos), entonces $X\simeq Y$. (El regreso no es cierto en general).
    \end{ejercicio}

\begin{ejercicio}
    % ID: 5
        Mostrar que un espacio $X$ es contraible si y sólo si existe una homotopía $H:X\times I\to X$ tal que $H(-,0)=\id_X$ y $H(-,1)=c_{x_0}$, para algún $x_0\in X$.
    \end{ejercicio}

\begin{ejercicio}
    % ID: 6
        Ver que $\R^n$ y $D^n$ son contraibles.
    \end{ejercicio}

\begin{ejercicio}
    % ID: 7
    \label{eje:Pi_1_es_invariante}
        Sean 
        $
            \begin{tikzcd}[column sep=0.5cm]
            (X,x_0) \arrow[r, "f"] & (Y,y_0)\arrow[r, "g"] & (Z,z_0)
        \end{tikzcd}
        $
        funciones continuas basadas. Mostrar que
        \begin{enumerate}
            \item [(1)] ${(g\circ f)}_{\#}=g_{\#}\circ f_{\#}$ y
            \item[(2)] ${\id_{(X,x_0)}}_{\#}=\id_{\Pi_1(X,x_0)}.$ 
        \end{enumerate}
    \end{ejercicio}

\begin{ejercicio}
    % ID: 8
        Verificar que la multiplicación por escalar y la suma en $\R^n$ son continuas.
    \end{ejercicio}

\begin{ejercicio}
    % ID: 9
        Mostrar que $\overline{\alpha}*\alpha\simeq c_{x_1}\rel$.
    \end{ejercicio}

\begin{ejercicio}
    % ID: 10
        Mostrar que $\Vec{\sigma}$ es un homomorfismo de grupo.
    \end{ejercicio}

\begin{ejercicio}
    % ID: 11
        Mostrar que $\Vec{\sigma}$ satisface las siguientes propiedades
        \begin{enumerate}
            \item[(1)] Si $\sigma\simeq\tau\rel$, entonces $\Vec{\sigma}=\Vec{\tau}.$
            \item[(2)] $\Vec{c}_{x_0}=\id_{\Pi_1(X,x_0)}.$ 
            \item[(3)]  Sea $\sigma_0:I\to X$ una trayectoria de $x_0$ a $x_1$ y $\sigma_1$ una trayectoria de $x_1$ a $x_2$. Mostrar que $$\overrightarrow{\sigma_0*\sigma_1}=\Vec{\sigma}_0\circ \Vec{\sigma}_1.$$
        \end{enumerate}
    \end{ejercicio}

\begin{ejercicio}
    % ID: 12
    Mostrar que $\p$ es la inversa de $p^{(X,x_0)
    }$.
\end{ejercicio}

\begin{ejercicio}
    % ID: 13
    Sea $F:\C\to\D$ un funtor. Mostrar que $F$ manda objetos isomorfos en $\C$ a objetos isomorfos en $\D$.
\end{ejercicio}

\begin{ejercicio}
    % ID: 14
        Mostrar que $\HTop$ es una categoría y describir qué es un isomorfismo en está categoría.
    \end{ejercicio}

\begin{ejercicio}
    % ID: 15
    Para cada espacio $(X,x_0)$, hagamos 
    \begin{align*}
        T_{(X,x_0)}:[(I^q/\partial I^q,*),(X,x_0)]&\to[(S^q,*),(X,x_0)]\\
        [\alpha]&\mapsto [\alpha\circ h],
    \end{align*}
    donde $h$ es el homeomorfismo presentado en \ref{eq:iso_In/partial_In_to_Sn}. Mostrar que $T$ es una equivalencia natural, i.e., para cualquier función continua basada $f:(X,x_0)\to(Y,y_0)$, el siguiente diagrama conmuta 
     \[
    \begin{tikzcd}[column sep=1.2cm, row sep=1.2cm] 
        \bIq\arrow[r,"\px"] \arrow[d,"f_+"'] & \text{[}(S^q,*),(X,x_0)\text{]}
        \arrow[d,"f_\#"']\\
        \YIq\arrow[r,"\py"] &\text{[}(S^q,*),(Y,y_0){]}
\end{tikzcd}
    \]
\end{ejercicio}
